%%
%% This is file `sample-acmsmall.tex',
%% generated with the docstrip utility.
%%
%% The original source files were:
%%
%% samples.dtx  (with options: `all,journal,bibtex,acmsmall')
%% 
%% IMPORTANT NOTICE:
%% 
%% For the copyright see the source file.
%% 
%% Any modified versions of this file must be renamed
%% with new filenames distinct from sample-acmsmall.tex.
%% 
%% For distribution of the original source see the terms
%% for copying and modification in the file samples.dtx.
%% 
%% This generated file may be distributed as long as the
%% original source files, as listed above, are part of the
%% same distribution. (The sources need not necessarily be
%% in the same archive or directory.)
%%
%%
%% Commands for TeXCount
%TC:macro \cite [option:text,text]
%TC:macro \citep [option:text,text]
%TC:macro \citet [option:text,text]
%TC:envir table 0 1
%TC:envir table* 0 1
%TC:envir tabular [ignore] word
%TC:envir displaymath 0 word
%TC:envir math 0 word
%TC:envir comment 0 0
%%
%% The first command in your LaTeX source must be the \documentclass
%% command.
%%
%% For submission and review of your manuscript please change the
%% command to \documentclass[manuscript, screen, review]{acmart}.
%%
%% When submitting camera ready or to TAPS, please change the command
%% to \documentclass[sigconf]{acmart} or whichever template is required
%% for your publication.
%%
%%
\documentclass[acmsmall]{acmart}

%% Authors custom
\usepackage{graphicx} % Required for inserting images
\usepackage{hyperref}
\usepackage{subcaption}

%%
%% \BibTeX command to typeset BibTeX logo in the docs
\AtBeginDocument{%
  \providecommand\BibTeX{{%
    Bib\TeX}}}

%% Rights management information.  This information is sent to you
%% when you complete the rights form.  These commands have SAMPLE
%% values in them; it is your responsibility as an author to replace
%% the commands and values with those provided to you when you
%% complete the rights form.
\setcopyright{acmlicensed}
\copyrightyear{2026}
\acmYear{2026}
\acmDOI{XXXXXXX.XXXXXXX}

%%
%% These commands are for a JOURNAL article.
\acmJournal{JACM}
\acmVolume{37}
\acmNumber{4}
\acmArticle{111}
\acmMonth{8}

%%
%% Submission ID.
%% Use this when submitting an article to a sponsored event. You'll
%% receive a unique submission ID from the organizers
%% of the event, and this ID should be used as the parameter to this command.
%%\acmSubmissionID{123-A56-BU3}

%%
%% For managing citations, it is recommended to use bibliography
%% files in BibTeX format.
%%
%% You can then either use BibTeX with the ACM-Reference-Format style,
%% or BibLaTeX with the acmnumeric or acmauthoryear sytles, that include
%% support for advanced citation of software artefact from the
%% biblatex-software package, also separately available on CTAN.
%%
%% Look at the sample-*-biblatex.tex files for templates showcasing
%% the biblatex styles.
%%

%%
%% The majority of ACM publications use numbered citations and
%% references.  The command \citestyle{authoryear} switches to the
%% "author year" style.
%%
%% If you are preparing content for an event
%% sponsored by ACM SIGGRAPH, you must use the "author year" style of
%% citations and references.
%% Uncommenting
%% the next command will enable that style.
%%\citestyle{acmauthoryear}

\newcommand{\hugo}[1]{\textcolor{violet}{#1}}
\newcommand{\chema}[1]{\textcolor{red}{#1}}
\newcommand{\ignacio}[1]{\textcolor{blue}{#1}}

%%
%% end of the preamble, start of the body of the document source.
\begin{document}

%%
%% The "title" command has an optional parameter,
%% allowing the author to define a "short title" to be used in page headers.
\title{Where Does Your Data Really Go? Rethinking Geographic Control in Anycast Systems}

%%
%% The "author" command and its associated commands are used to define
%% the authors and their affiliations.
%% Of note is the shared affiliation of the first two authors, and the
%% "authornote" and "authornotemark" commands
%% used to denote shared contribution to the research.
\author{Hugo Pascual}
\orcid{0009-0003-8357-9037}
\email{hugo.pascual.adan@upm.es}
\affiliation{
  \institution{Information Processing and Telecommunications Center, Universidad Politécnica de Madrid}
  \department{ETSI Telecomunicación}
  \city{Madrid}
  \country{Spain}
}

\author{Ignacio Castro}
\orcid{0000-0002-7739-6184}
\email{i.castro@qmul.ac.uk}
\affiliation{
  \institution{Queen Mary University of London}
  \city{London}
  \country{United Kingdom}
}

\author{Jose M. del Alamo}
\orcid{0000-0002-6513-0303}
\email{jm.delalamo@upm.es}
\affiliation{
  \institution{Information Processing and Telecommunications Center, Universidad Politécnica de Madrid}
  \department{ETSI Telecomunicación}
  \city{Madrid}
  \country{Spain}
}

%%
%% By default, the full list of authors will be used in the page
%% headers. Often, this list is too long, and will overlap
%% other information printed in the page headers. This command allows
%% the author to define a more concise list
%% of authors' names for this purpose.
\renewcommand{\shortauthors}{Pascual, Castro and Del Alamo}

%%
%% The abstract is a short summary of the work to be presented in the
%% article.
\begin{abstract}
Anycast routing is widely used by major web services to achieve global reach, load balancing, and resilience against denial-of-service attacks. However, the same mechanism that optimizes performance can complicate personal data protection and data sovereignty requirements. Anycast design prioritizes network efficiency over geographic determinism: user traffic is directed to the topologically nearest server, which may reside outside the user’s legal or jurisdictional region. This behavior challenges assumptions underlying modern personal data protection and data sovereignty frameworks, which often require control over where data is processed. Our measurements across a globally distributed set of anycast deployments show that cross-border routing is not an edge case but a systemic property of today’s Internet. In practice, sensitive data frequently crosses jurisdictional boundaries, often without service providers being explicitly aware. To address this gap, we argue that deterministic geographic control requires an architectural approach that moves jurisdictional policy to the network edge rather than modifying Internet routing. We demonstrate this approach through a simple, deployable prototype that combines DNS-based regionalization with controlled anycast deployments. This approach enables service providers to enforce explicit geographic constraints on data processing while preserving the operational benefits of anycast. Our results suggest that the tension between performance and jurisdictional control is not inherent to anycast, but to deployments that leave geographic constraints implicit.
\end{abstract}

%%
%% The code below is generated by the tool at http://dl.acm.org/ccs.cfm.
%% Please copy and paste the code instead of the example below.
%%
\begin{CCSXML}
<ccs2012>
   <concept>
       <concept_id>10002978.10003029.10011150</concept_id>
       <concept_desc>Security and privacy~Privacy protections</concept_desc>
       <concept_significance>500</concept_significance>
       </concept>
   <concept>
       <concept_id>10003033.10003083.10011739</concept_id>
       <concept_desc>Networks~Network privacy and anonymity</concept_desc>
       <concept_significance>500</concept_significance>
       </concept>
   <concept>
       <concept_id>10003033.10003034.10003035.10003037</concept_id>
       <concept_desc>Networks~Naming and addressing</concept_desc>
       <concept_significance>300</concept_significance>
       </concept>
 </ccs2012>
\end{CCSXML}

\ccsdesc[500]{Security and privacy~Privacy protections}
\ccsdesc[500]{Networks~Network privacy and anonymity}
\ccsdesc[300]{Networks~Naming and addressing}

%%
%% Keywords. The author(s) should pick words that accurately describe
%% the work being presented. Separate the keywords with commas.
\keywords{Anycast, DNS, Data Protection, Data Sovereignty, Privacy}

\received{8 August 2026}
\received[revised]{1 Month 2026}
\received[accepted]{1 Month 2026}

%%
%% This command processes the author and affiliation and title
%% information and builds the first part of the formatted document.
\maketitle

\section{Where Does Your Data Really Go?}
\label{sec:introduction}
You have just uploaded laboratory results related to a recent medical diagnosis. The information contained describes your symptoms, and discloses your medical history. You have entered this deeply personal information from your own home, using an app available and marketed in your country. You would reasonably expect this data to benefit from the legal protections that apply there. The underlying network may, however, route the communication to infrastructure located in another country and jurisdiction. Our measurement campaign confirms that this is not merely hypothetical: traffic originating in the United States was observed being routed to infrastructure in South Africa, Japan, or Cuba, among others, despite the available infrastructure in the United States. This happens because Internet routing optimizes performance rather than jurisdictional boundaries. Unfortunately, this behavior is usually invisible to users and, in many cases, even to service providers; yet it occurs far more frequently than commonly assumed. As digital services become increasingly central, understanding and controlling where sensitive data is processed is a fundamental technical and societal challenge.

Data sovereignty is a central concern for organizations operating digital services. Regulations such as the General Data Protection Regulation (GDPR) in Europe, together with similar legal frameworks in other jurisdictions, require organizations to ensure that personal data is processed within approved legal boundaries or transferred only under well-defined legal safeguards. While these obligations are expressed in legal terms, satisfying them ultimately depends on the technical mechanisms that determine where Internet communications are processed.

One of the technologies that underpins modern Internet services is anycast. CDNs, cloud platforms, authoritative DNS infrastructures, and DDoS protection services rely on it to provide low latency, scalability, and resilience. Multiple geographically distributed servers advertise the same IP address, allowing the network to transparently select a service instance. However, this selection is driven by network routing rather than jurisdictional considerations, meaning that neither users nor providers can reliably predict where an anycast communication will ultimately be processed.

Because routing decisions are determined by the Border Gateway Protocol (BGP) and the independent policies of Internet Service Providers, communications may cross jurisdictional boundaries even when suitable infrastructure exists within the user's own region. Previous work has shown that Internet routing frequently diverges from geographic intuition, producing international routing detours and unexpected cross-border paths~\cite{Comarela2016DetectingApplications, Shah2016TowardsDetours, Sun2021SecuringAttacks}. This creates a fundamental mismatch between legal obligations and technical capabilities: organizations are increasingly required to control where data is processed, while cloud platforms and CDNs deliberately abstract away routing decisions to maximize performance and availability.

Addressing this challenge requires answering two complementary questions: (1) can service providers determine where communications are processed? and, (2) can they ensure that traffic remains within a designated jurisdiction? Existing approaches—including modified anycast deployments, BGP-based routing strategies, DNS server selection, and IP geolocation—primarily optimize routing or performance rather than enforce deterministic jurisdictional boundaries~\cite{Zhu2022TheWorlds, Cozar2022ReliabilityTransfers}.

A natural question is whether these limitations could be addressed by modifying Internet routing itself. In practice, this is neither realistic nor desirable. Routing protocols such as BGP are designed to provide scalable connectivity across independently administered networks rather than enforcing application-specific policies. Following the long-standing Internet design principle of moving application-specific policy to the edge, we argue that jurisdictional constraints should likewise be enforced before traffic enters the anycast infrastructure. We demonstrate this approach through Region-Constrained Anycast over DNS (RCA-DNS), a simple prototype that performs jurisdiction selection before DNS resolution, decoupling policy decisions from routing enforcement while preserving conventional anycast's scalability, resilience, and performance.

\section{How Anycast Works in Practice}
\label{sec:background}

Anycast has become a fundamental building block of modern Internet services, including CDNs, authoritative DNS infrastructures, cloud platforms, and DDoS protection services. Multiple geographically distributed servers advertise the same IP prefix, allowing clients to communicate with what appears to be a single destination, while the network transparently selects one of the available service instances.

Although anycast is often described as routing traffic to the "nearest" server, this description is only an approximation. On the public Internet, anycast routing is primarily determined by the Border Gateway Protocol (BGP), the inter-domain routing protocol used by Internet Service Providers (ISPs). Each anycast instance advertises the same IP prefix, and every ISP independently selects which advertisement to follow according to its local BGP policies. These policies are driven primarily by routing preferences, commercial agreements, and traffic engineering considerations rather than by geographic distance. Consequently, the selected server is not necessarily the geographically closest one, but the one that is most preferred according to the distributed BGP decision process.

This distinction has important implications beyond performance. Because BGP has no notion of political borders, legal jurisdictions, or data sovereignty requirements, routing decisions may direct traffic across countries or regulatory regions even when local infrastructure exists. Previous studies have documented international routing detours and unexpected paths caused by BGP policies, showing that Internet routing frequently diverges from geographic intuition \cite{Comarela2016DetectingApplications, Shah2016TowardsDetours, Sun2021SecuringAttacks}. These observations demonstrate that routing policies can inadvertently violate jurisdictional expectations despite operating correctly from a networking perspective.

A substantial body of research has examined anycast from the perspectives of performance, scalability, and operational reliability. Previous work has evaluated global deployments across critical Internet infrastructures, such as root DNS services and large-scale CDNs, considering metrics including latency, routing stability, replica selection, and user-to-server mapping dynamics \cite{Li2023EvaluatingScope, Calder2015AnalyzingCDN, Koch2021AnycastSystems}. Other studies have investigated how routing policies, infrastructure placement, and peering arrangements influence the regional performance of anycast services \cite{Zhou2023RegionalPotentials}.

However, these approaches do not address the fundamental challenge of enforcing where traffic is processed when geographic or jurisdictional boundaries matter. Their primary objectives are routing efficiency and performance rather than deterministic jurisdictional guarantees.

Zhu et al.~\cite{Zhu2022TheWorlds} propose two techniques combining unicast, anycast, and modified BGP announcements to increase operator control over anycast deployments. While their approach improves routing control under specific scenarios, it requires significant network modifications and still cannot deterministically constrain where traffic is processed. Even their fastest failover configuration forwards approximately 82\% of requests to unintended service instance jurisdictions.

Jia et al.~\cite{Jia2019DeployingInternet} combine DNS-based replica selection with active probing to improve server selection and reduce connection latency, while Evang et al.~\cite{Evang2024AnycastTuning} define performance metrics and BGP tuning techniques to improve client-to-server mapping, latency, and load distribution. Although valuable for operational optimization, neither approach provides deterministic guarantees that communications remain within a designated legal jurisdiction.

The absence of deterministic jurisdictional guarantees has become increasingly important as governments and regulatory bodies have introduced explicit constraints on international data transfers. Regulations such as the General Data Protection Regulation (GDPR) in the European Economic Area (EEA) \cite{REGULATIONSRelevance}, the UK GDPR and Data Protection Act \cite{Data2018}, and similar legal frameworks in the United States, Australia, and China \cite{Rep.Pallone2024Text2024, ReadOAIC, TheTranslation} require organizations to ensure that personal data is transferred in compliance with specific legal safeguards and, in many cases, to disclose the jurisdictions in which it is processed.

Meeting these requirements is particularly challenging for conventional anycast deployments. Since routing decisions emerge from independently administered BGP policies, service providers cannot deterministically control which anycast instance ultimately processes a given request. Consequently, even providers operating infrastructure in multiple jurisdictions cannot reliably guarantee that user traffic remains within the region required by their data governance policies. Recent empirical evidence confirms that this is not merely a theoretical concern. Pascual et al.~\cite{Pascual2025AnycastDisaster} showed that privacy policies frequently fail to accurately disclose the jurisdictions in which personal data may be processed because conventional anycast deployments lack mechanisms for deterministically constraining the geographic destination of communications.

Despite these challenges, abandoning anycast is rarely a viable solution. Anycast remains a cornerstone of modern Internet infrastructures because it provides low latency, resilience against denial-of-service attacks, fault tolerance, and efficient load distribution. The challenge is therefore not to replace anycast but to complement it with mechanisms that preserve these operational advantages while providing deterministic jurisdictional guarantees.

\section{Putting boundaries to regions with anycast}
\label{sec:solution}

We present Region-Constrained Anycast over DNS (RCA-DNS), a simple and deployable approach that enables service providers to retain the benefits of anycast while enforcing geographic constraints on where user data is processed. 

The key idea is to decouple two concerns that are traditionally intertwined: determining a user’s region and controlling where traffic is routed. RCA-DNS separates these functions by combining region-specific DNS records with controlled anycast deployments, introducing explicit geographic boundaries into otherwise unconstrained anycast systems.

\subsection{Region-Constrained Anycast over DNS}

The approach defines multiple DNS records for the same service, each corresponding to a distinct geographic region. Each record resolves to a different anycast IP address, and each IP is announced exclusively by servers located within the corresponding region. As a result, traffic directed to a given DNS record is confined to infrastructure within the predefined geographic boundary, regardless of how the user’s location is determined.

For example, consider a user in Finland whose traffic might otherwise be routed to infrastructure in Russia. If the service instead directs the user to \textit{eea.application.com}, the request is guaranteed to be served by infrastructure within the European Economic Area (EEA). The same principle can be extended to other regions (e.g., \textit{usa.application.com}), allowing providers to enforce jurisdiction-specific constraints \cite{Rep.Pallone2024Text2024, ECFRAdversaries.}. Importantly, the definition of a region is flexible and can be adapted to specific legal or organizational requirements. For instance, a provider could define an \textit{eea-extended.application.com} region that includes the EEA together with the United Kingdom and Switzerland, reflecting common data-transfer arrangements and regulatory practices in Europe. 

At the source of the communication (e.g., a mobile application), the service selects the appropriate DNS record based on the region associated with the user or the traffic. This decision may rely on user-declared information, coarse geolocation, or other policy-relevant attributes. Once the DNS resolution is completed, traffic delivery relies entirely on standard anycast routing, preserving scalability and performance while ensuring regional confinement.

The mechanism intentionally separates region determination from routing control. RCA-DNS only requires that a jurisdiction be selected before DNS resolution; how that jurisdiction is determined is outside the scope of the routing architecture. Service providers may rely on user-provided information, coarse IP geolocation, contractual relationships, enterprise policies, regulatory requirements, or other application-specific mechanisms. Once the jurisdiction has been selected, RCA-DNS guarantees that traffic remains confined to the corresponding region.

\subsection{Deployment Considerations and Summary}

RCA-DNS requires one anycast IP address per region-specific DNS record. In practice, this is manageable: IPv6 provides ample address space, while IPv4 addresses can be shared across services, as is commonly done by infrastructure providers. Traffic can then be internally demultiplexed using higher-layer information, such as the destination domain name, while ensuring that each address is announced only within its intended region.

The architecture is fully compatible with today's Internet infrastructure. Region-specific DNS records are managed within authoritative zones under the provider's control and do not increase the load on DNS root or top-level infrastructure. By combining DNS-based regionalization with controlled anycast announcements, RCA-DNS provides deterministic control over where traffic is processed while preserving the scalability, performance, and resilience of conventional anycast.

\section{Real world measurements}
\label{sec:experiment}

We empirically evaluate both the prevalence of cross-border routing in current anycast deployments and the practical impact of enforcing geographic constraints.

The first experiment quantifies the extent of cross-border routing in operational anycast infrastructures, while the second evaluates the latency implications of introducing deterministic regional boundaries.

To support reproducibility, the analysis code and datasets used in this study are publicly available\footnote{Global Cross Border Analysis Code: \url{https://github.com/STRAST-UPM/experiments_hunter}. Global Cross Border Data: \url{https://doi.org/10.17632/n93r278vzy.1}. RCA-DNS evaluation code: \url{https://github.com/STRAST-UPM/RCA-DNS}. RCA-DNS evaluation data: \url{https://doi.org/10.17632/ntys787cys.1}}.

\subsection{Global Measurement of Cross-Border Anycast Routing}

\subsubsection{Methodology}

To quantify how often anycast routing results in cross-border data transfers, we conducted a large-scale measurement campaign using \emph{Hunter}, an anycast measurement framework \cite{Pascual2025AnycastDisaster}.

We analyzed 195 anycast IP addresses associated with Android applications that transmit personal data. These IPs represent service endpoints observed in real-world deployments rather than a comprehensive sample of all anycast deployments. The applications were selected from popular Google Play categories and prioritized by download count. For each IP address, we conducted ten measurement rounds using RIPE Atlas probes distributed across 183 countries, randomly selecting ten probes per round. Measurements were collected during two six-day periods in July and September 2025 at irregular intervals to capture routing variability.

\subsubsection{Results}

After filtering out incomplete traces, non-responsive paths, and applying a criterion of three origin-destination pairs for successful observations for a valid result, the dataset contained 504,690 valid routing tracks. Each track represents the observed path between a probe location and a specific instance of an anycast address.

\begin{figure}[t]
    \centering
    \begin{subfigure}{0.48\textwidth}
        \centering
        \begin{tabular}{lr}
            \hline
                Metric & Value \\
            \hline
                Applications analyzed & 379 \\
                Anycast IP addresses & 195 \\
                Countries covered & 183 \\
                Valid routing tracks & 504,690 \\
                Cross-border tracks & 58.36\% \\
                IPs with cross-border routing & 99.49\% \\
            \hline
        \end{tabular}
        \caption{Metrics summary table.}
        \label{tab:global_hunter_summary}
    \end{subfigure}
    \hfill
    \begin{subfigure}{0.48\textwidth}
        \centering
        \includegraphics[width=\textwidth]{figure_tracks_destinations_color_from_usa_no_scale.png}
        \caption{Destinations for traffic originated in USA.}
        \label{fig:destinations_from_usa}
    \end{subfigure}

    \vspace{0.5em}

    \begin{subfigure}{0.48\textwidth}
        \centering
        \includegraphics[width=\textwidth]{figure_tracks_destinations_color_from_eea_no_scale.png}
        \caption{Destinations for traffic originated in EEA.}
        \label{fig:destinations_from_eea}
    \end{subfigure}
    \hfill
    \begin{subfigure}{0.48\textwidth}
        \centering
        \includegraphics[width=\textwidth]{figure_tracks_destinations_color_from_japan_no_scale.png}
        \caption{Destinations for traffic originated in Japan.}
        \label{fig:destinations_from_japan}
    \end{subfigure}

    \caption{Visualization and summary of the global anycast cross-border measurement campaign.}
    \label{fig:visualization_global_hunter}
\end{figure}

\paragraph{Prevalence of cross-border routing}

Cross-border routing was widespread. Overall, 58.36\% of valid tracks terminated in a country different from the origin, and 99.49\% of the analyzed anycast IPs exhibited at least one cross-border routing event. As summarized in Table~\ref{tab:global_hunter_summary}, cross-border routing appears to be the norm in contemporary anycast deployments rather than an exceptional behavior.

\paragraph{Regional routing patterns}

Several regional patterns emerged. Traffic originating in Latin America frequently terminated in the USA, highlighting the central role of US infrastructure. Conversely, US-originated traffic (Figure~\ref{fig:destinations_from_usa}) that crossed the border was also observed terminating in countries such as Brazil, Japan, South Africa, the EEA, and, in some cases, Cuba, a country designated as a \textit{foreign adversary} by U.S. regulation \cite{ECFRAdversaries.}.

The EEA exhibited substantial external routing, with 34.88\% of communications terminating outside the region, including destinations such as Russia, India, and South Africa (Figure~\ref{fig:destinations_from_eea}), which do not provide a level of data protection equivalent to that of the EU. In contrast, Japan (Figure~\ref{fig:destinations_from_japan}), India, and South Korea showed stronger regional retention, while African traffic frequently terminated in Europe, reflecting continued dependence on external infrastructure.

\subsection{Performance Impact of Region-Constrained Anycast}

\subsubsection{Methodology}\label{sss:Methodology}

The second experiment evaluates the performance impact of enforcing jurisdictional constraints through the RCA-DNS prototype. To this end, we compare two deployment strategies built on the same underlying infrastructure: a conventional global anycast deployment and the proposed region-constrained deployment.

We deployed an identical HTTP service in every available Google Cloud region, excluding the Middle East regions for technical reasons. All instances simultaneously advertised the domain \textit{global.anycastprivacy.org}, emulating a conventional global anycast service. In parallel, each jurisdiction was assigned a dedicated regional domain (e.g., \textit{us.anycastprivacy.org}, \textit{eea-expanded.anycastprivacy.org}), whose anycast address was announced exclusively by the instances located within the corresponding region. This configuration implements the RCA-DNS architecture while keeping the deployed infrastructure identical across both approaches. 

Google Cloud defines eight geographic regions with different numbers of deployment locations, ranging from a single location in \textit{Africa} to thirteen across \textit{EEA-Extended} (eleven within the EEA, one in the United Kingdom, and one in Switzerland). The remaining regions are \textit{Asia} (9 locations), \textit{the United States} (9), \textit{North America} (Canada and Mexico, 3), \textit{Australia} (2), and \textit{South America} (2). In our deployment, each region is exposed through a dedicated regional domain (e.g., \textit{eea-extended.anycastprivacy.org}), while all locations simultaneously participate in the conventional global deployment.

\begin{comment}
\begin{table}[ht]
    \centering
    \begin{tabular}{lccc}
        \toprule
        \textbf{Domain name} & \textbf{Number of regions} \\
        \midrule
        africa.anycastprivacy.org       & 1 \\
        asia.anycastprivacy.org         & 9 \\
        eea-extended.anycastprivacy.org & 13 \\
        northamerica.anycastprivacy.org & 3 \\
        australia.anycastprivacy.org    & 2 \\
        southamerica.anycastprivacy.org & 2 \\
        us.anycastprivacy.org           & 9 \\
        \bottomrule
    \end{tabular}
    \caption{Number of Google Cloud regions for domain name.}
    \label{tab:gc_regions_per_domain}
\end{table}
\end{comment}

Latency measurements were performed using RIPE Atlas probes. Up to 10 probes were randomly selected from every country worldwide, subject to probe availability. Measurements were conducted between July 21 and July 28, 2026, at six-hour intervals. During each measurement round, probes issued HTTP requests to both the global and every regional service, and the round-trip time (RTT) was recorded. Probes were selected worldwide to evaluate both the impact on users located within the protected region and the performance implications for users located outside it. RTT was selected as the primary performance metric because it captures both network propagation delay and server responsiveness.

\subsubsection{Results}

\begin{table}[ht]
    \centering
    \begin{tabular}{lccc}
        \toprule
        \textbf{Domain} & \textbf{Regional RTT (ms)} & \textbf{Global RTT (ms))} & \textbf{Change} \\
        \midrule
        africa.anycastprivacy.org & 369.86 & 143.34 & -158.03\% \\
        asia.anycastprivacy.org & 339.25 & 128.51 & -163.99\% \\
        eea-extended.anycastprivacy.org & 48.43 & 51.25 & +5.50\% \\
        northamerica.anycastprivacy.org & 104.77 & 50.38 & -107.96\% \\
        australia.anycastprivacy.org & 125.27 & 120.99 & -3.54\% \\
        southamerica.anycastprivacy.org & 772.05 & 118.97 & -548.93\% \\
        us.anycastprivacy.org & 153.04 & 145.31 & -5.32\% \\
        \bottomrule
    \end{tabular}
    \caption{Median RTT observed from probes within each protected region to regional and global anycast deployments.}
    \label{tab:median_rtt_comparison_regions}
\end{table}

\begin{comment}
\begin{table}[t]
    \centering
    \begin{tabular}{lccc}
        \toprule
        \textbf{Region} & \textbf{Regional RTT (ms)} & \textbf{Global RTT (ms)} & \textbf{Change} \\
        \midrule
        africa.anycastprivacy.org & 454.30 & 191.82 & -136.84\% \\
        asia.anycastprivacy.org & 507.95 & 224.36 & -126.39\% \\
        eea-extended.anycastprivacy & 61.05 & 189.81 & +67.85\% \\
        northamerica.anycastprivacy.org & 371.96 & 213.08 & -74.56\% \\
        australia.anycastprivacy.org & 341.65 & 204.62 & -66.97\% \\
        southamerica.anycastprivacy.org & 586.87 & 156.75 & -274.53\% \\
        us.anycastprivacy.org & 1150.78 & 329.70 & -249.04\% \\
        \bottomrule
    \end{tabular}
    \caption{Average RTT observed from probes within each protected region to regional and global anycast deployments.}
    \label{tab:average_rtt_comparison_regions}
\end{table}
\end{comment}

Overall, the results show that the performance impact of RCA-DNS is not uniform across jurisdictions, as shown in Table~\ref{tab:median_rtt_comparison_regions}. Instead, it depends primarily on the availability and geographic distribution of infrastructure within the protected region. Regions with a sufficiently dense regional deployment can enforce jurisdictional boundaries with little or no latency penalty, whereas regions with sparse infrastructure experience higher RTTs because traffic can no longer benefit from nearby infrastructure located outside the protected jurisdiction or from routing paths optimized through neighboring regions. This relationship is consistent with the distribution of Google Cloud regions summarized in section ~\ref{sss:Methodology}.

\begin{figure}[ht]
    \begin{subfigure}{0.48\textwidth}
        \includegraphics[width=\textwidth]{figure_cdf_rtts_from_eea-extended_to_eea-extended.anycastprivacy.org.png}
        \caption{RTTs CDF observed from EEA-Extended countries to anycast instances in the domain eea-extended.anycastprivacy.org}
        \label{fig:cdf_rtt_eea-extended_to_eea-extended_rca-dns}
    \end{subfigure}
    \hspace{0.02\textwidth}
    \begin{subfigure}{0.48\textwidth}
        \includegraphics[width=\textwidth]{figure_cdf_rtts_from_eea-extended_to_global.anycastprivacy.org.png}
        \caption{RTTs CDF observed from EEA-Extended countries to anycast instances in the domain global.anycastprivacy.org}
        \label{fig:cdf_rtt_eea-extended_to_global_rca-dns} 
    \end{subfigure}
    \caption{Comparison of RTT distributions observed by EEA-Extended probes when accessing the service through the proposed regional deployment (left) and a conventional global anycast deployment (right).}
\end{figure}

EEA-Extended provides a representative example of a well-provisioned jurisdiction. As shown in Figures~\ref{fig:cdf_rtt_eea-extended_to_eea-extended_rca-dns} and~\ref{fig:cdf_rtt_eea-extended_to_global_rca-dns}, the regional deployment not only preserves performance but also reduces the occurrence of high-latency requests. These probes observe a 5.5\% reduction in median RTT (48.43~ms versus 51.25~ms), while the average RTT improves by 67.85\% (61.05~ms versus 189.81~ms). These results suggest that constraining traffic to infrastructure within the protected region can eliminate routing detours that occur under conventional global anycast while maintaining the responsiveness expected from modern Internet services.

The remaining regions exhibit behavior consistent with their available infrastructure. Jurisdictions with relatively dense deployments, such as the United States, experience only modest median RTT increases of 5.3\%. In contrast, Africa and South America incur substantially larger penalties because their limited regional infrastructure forces many users to communicate with geographically distant service instances. These observations indicate that the performance cost of regional jurisdictional enforcement is determined primarily by infrastructure availability rather than by the prototype itself. Consequently, as cloud providers and CDNs continue expanding their regional footprints, the latency overhead associated with jurisdiction-constrained anycast is expected to decrease.

\section{Implications for stakeholders and limitations}
\label{sec:discussion}

Our results show that cross-border routing in anycast infrastructures is widespread, as nearly all analyzed anycast IP addresses exhibited at least one cross-border routing event. International data transfers are not inherently unlawful. Many regulatory frameworks allow them under specific legal mechanisms. However, the prevalence and dynamic nature of the routing patterns observed in this study raise important questions regarding transparency, accountability, and effective control. When routing decisions are delegated entirely to global anycast optimization mechanisms, service providers may lack precise knowledge of where data is processed at any given moment.

This uncertainty complicates compliance with modern personal data protection and data sovereignty requirements. Even when service providers intend to restrict processing to specific jurisdictions, the absence of deterministic geographic constraints makes it difficult to ensure that traffic consistently remains within defined legal boundaries.

Our second experiment shows that regional enforcement through DNS-based routing is technically feasible and does not impose prohibitive performance penalties within the protected region. Latency increases primarily affect users located outside the constrained region, while users inside the region experience performance comparable to conventional CDN-backed deployments.

\paragraph{Implications for Service and Infrastructure Providers}

Our results suggest that neither service providers nor infrastructure operators should assume that conventional anycast deployments respect jurisdictional boundaries. When geographic constraints matter, service providers require mechanisms that provide deterministic control over where communications are processed rather than relying on routing behavior that emerges from independent network policies.

Our prototype demonstrates that such control can be introduced without modifying Internet routing protocols. More broadly, this suggests that jurisdictional policy can be separated from routing enforcement while preserving the operational advantages of anycast.

These findings also have implications for cloud and CDN providers. Greater transparency regarding traffic locality and first-class support for region-scoped anycast services would allow customers to better align technical deployments with legal and regulatory obligations. As data sovereignty becomes an increasingly important consideration in digital infrastructure policy, explicit mechanisms for geographic service scoping may become a key capability of future Internet platforms.

\paragraph{User Location and Jurisdiction Determination}

RCA-DNS intentionally decouples jurisdiction determination from routing enforcement. The proposed architecture only requires that the applicable jurisdiction be determined before DNS resolution; it does not prescribe how this decision should be made. This separation allows the routing mechanism to remain independent of application-specific policy decisions while enabling deterministic enforcement of jurisdictional boundaries.

Depending on the application, the applicable jurisdiction may be determined through authenticated user profiles, contractual information, enterprise policies, coarse IP geolocation, device-level location signals, or explicit user declarations. Each mechanism presents different trade-offs in terms of accuracy, deployability, and privacy, making a single solution inappropriate for all application domains.

Future advances in jurisdiction determination can therefore be incorporated into RCA-DNS without requiring changes to the routing architecture since routing enforcement remains independent of the policy mechanism used to assign users to jurisdictions.

\paragraph{Limitations}

This study focuses on routing-level observations and does not assess application-layer processing, storage practices, or contractual safeguards that may also influence regulatory compliance. Measurement visibility also varies across countries, as illustrated by the case of China, which limits interpretability in certain jurisdictions.

In addition, the latency evaluation reflects controlled deployments and may differ in large-scale production environments with more extensive infrastructure distribution.

Despite these limitations, the combined evidence from both experiments indicates that international routing in anycast systems is widespread and that introducing deterministic regional constraints is technically achievable with a bounded performance impact.

\section{Conclusion}
\label{sec:conclusion}

Anycast has become a fundamental component of modern Internet infrastructures, but it was designed to optimize connectivity rather than to enforce jurisdictional boundaries. Consequently, the geographic location in which communications are ultimately processed often remains outside the control of both users and service providers.

Our measurements show that cross-border routing is widespread across contemporary anycast deployments, highlighting a mismatch between modern data governance requirements and the behavior of the underlying network infrastructure. These findings suggest that infrastructure location alone cannot provide reliable guarantees about where personal data is processed.

When deterministic geographic guarantees are required, we argue that the appropriate response is not to modify Internet routing, but to move jurisdictional policy to the network edge. Our RCA-DNS prototype demonstrates one way to realize this approach using existing DNS and anycast mechanisms, providing regional confinement without changes to routing protocols. More broadly, the results show that stronger geographic control can be introduced into today's Internet by separating jurisdictional policy from routing enforcement.

%%%%%%%%%

% \section{Declaration of generative AI and AI-assisted technologies in the writing process}

% During the preparation of this work, the authors used AI tools like Grammarly in order to improve the language and readability. After using this tool, the authors reviewed and edited the content as needed and take full responsibility for the content of the publication.

%%%%%%%%%

%%
%% The acknowledgments section is defined using the "acks" environment
%% (and NOT an unnumbered section). This ensures the proper
%% identification of the section in the article metadata, and the
%% consistent spelling of the heading.
\begin{acks}
Thanks to RIPE for granting access to the Atlas platform, which has been used to perform all the measurements needed for the experiments.
\end{acks}

%%
%% The next two lines define the bibliography style to be used, and
%% the bibliography file.
\bibliographystyle{ACM-Reference-Format}
\bibliography{references}

%% Use this line on submission to correct Proof View
%% Do not worry if the references does not appear in Overleaf
%% Obtain the file from the compiled files in Overleaf
%% \input{output.bbl}

%%
%% If your work has an appendix, this is the place to put it.
\appendix


\end{document}
\endinput
%%
%% End of file `sample-acmsmall.tex'.
